% Proposta de Projeto — Aprendizado de Máquina (UFPR)
% Baseado no template IEEE/CVPR

\documentclass[10pt,twocolumn,letterpaper]{article}

%%%%%%%%% PAPER TYPE
\usepackage{cvpr}

% Pacotes
\usepackage{graphicx}
\usepackage{amsmath}
\usepackage{amssymb}
\usepackage{booktabs}
\usepackage[utf8]{inputenc}
\usepackage[T1]{fontenc}
\usepackage[brazil]{babel}
\usepackage[pagebackref,breaklinks,colorlinks]{hyperref}
\usepackage[capitalize]{cleveref}
\crefname{section}{Seção}{Seções}
\Crefname{section}{Seção}{Seções}
\Crefname{table}{Tabela}{Tabelas}
\crefname{table}{Tab.}{Tabs.}

\def\cvprPaperID{*****}
\def\confName{CVPR}
\def\confYear{2022}

\begin{document}

%%%%%%%%% TÍTULO
\title{Proposta: Classificação de Eventos de Chuva Extrema
com Dados Meteorológicos do INMET}

\author{
Pedro Henrique Marques de Lima \\
Universidade Federal do Paraná \\
{\tt\small pedro.marques@ufpr.br}
\and
Felipe Pereira Gonçalves \\
Universidade Federal do Paraná \\
{\tt\small felipe.goncalves@ufpr.br}
}
\maketitle

% =====================================================================
\section{Problemas e Aplicações}

Enchentes e deslizamentos provocados por chuvas extremas causam
centenas de mortes e bilhões de reais em prejuízos no
Brasil~\cite{marengo2020trends,cenad2024}. Antecipar esses eventos a
partir de variáveis meteorológicas observadas é relevante para a
proteção civil.

\textbf{Problema~1 (tabular).}
Classificação binária horária: dado um vetor de variáveis
meteorológicas $\mathbf{x}\!\in\!\mathbb{R}^{17}$, prever se o
acumulado de precipitação excede um limiar crítico ($y\!=\!1$).
Os desafios são o desbalanceamento de classes (chuvas extremas são
raras), dados faltantes nas séries do INMET e a dependência temporal
entre observações.

\textbf{Problema~2 (não tabular).}


% =====================================================================
\section{Conjunto de Dados}

Os dados provêm do INMET\footnote{\url{https://portal.inmet.gov.br/dadoshistoricos}}
(observações horárias de estações automáticas, 2022--2025). Foram
selecionados cinco eventos reais de precipitação extrema
(\cref{tab:eventos}); para cada um, a estação mais próxima foi
escolhida entre 150 disponíveis nos estados alvo (PR, SP, RS, PE),
com recorte de $\pm$4~meses em torno do desastre.

\begin{table}[h]
\centering
\caption{Eventos selecionados e estações INMET.}
\label{tab:eventos}
\small
\setlength{\tabcolsep}{3pt}
\begin{tabular}{@{}llll@{}}
\toprule
\textbf{Região} & \textbf{UF} & \textbf{Estação} & \textbf{Data} \\
\midrule
Serra do Mar / Guaratuba  & PR & Morretes (A873)   & 28/11/2022 \\
Litoral Norte             & SP & Bertioga (A765)   & 19/02/2023 \\
Vale do Taquari           & RS & Teutônia (A882)   & 04/09/2023 \\
Bacia do Guaíba           & RS & P.\ Alegre (A801) & 02/05/2024 \\
Região Metropolitana      & PE & Palmares (A357)   & 06/02/2025 \\
\bottomrule
\end{tabular}
\end{table}

O conjunto consolidado possui \textbf{16\,013 registros horários} e
\textbf{17 variáveis}: precipitação, pressão atmosférica (3),
radiação global, temperaturas (5), umidade relativa (3) e vento~(3).
A radiação global apresenta $\sim$43\% de valores ausentes (período
noturno); demais lacunas são pontuais ($<$3\%), salvo a estação A873
($\sim$50\% em umidade/orvalho).

Inserir Conjunto de Dados do Problema 2 aqui.

% =====================================================================
\section{Metodologia}

\subsection{Pré-processamento}
\textbf{Problema 1:} Falhas serão tratadas por interpolação temporal
O rótulo $y\!=\!1$ será definido por limiar de precipitação
($\geq$20\,mm/h, calibrado pelos critérios INMET). Variáveis
derivadas de janelas deslizantes (acumulado 3--12\,h, variação de
pressão) serão adicionadas. Normalização por z-score. \\
\textbf{Problema 2:}

\subsection{Modelos}
\textbf{Problema 1:}
Random Forest, XGBoost/LightGBM.
Desbalanceamento tratado com SMOTE e ajuste de
pesos.
\textbf{Problema 2:}

\subsection{Avaliação}
\textbf{Problema 1:}
Métricas: F1-Score, AUC-ROC, Precision-Recall AUC e Balanced
Accuracy, com matrizes de confusão e curvas ROC/PR. Divisão
treino/validação/teste (70/15/15\%) respeitando ordem temporal.
Espera-se que \emph{ensembles} de árvores apresentem melhor
desempenho~\cite{grinsztajn2022tree}, com foco em maximizar
\emph{recall} da classe positiva.
\textbf{Problema 2:}

% =====================================================================
\newpage
{\small
\bibliographystyle{ieee_fullname}
\bibliography{egbib}
}

\end{document}
